% ======================================================================
\section{Deploying and testing for a Raspberry Pi 4}\label{sec:pi}
% ======================================================================

You do not own a Raspberry Pi. That is a normal situation and it does not stop
you from doing most of this properly --- provided you are precise about which
claims survive emulation and which do not. Conflating those two is how resource
claims get rejected in review.

\subsection{Three levels of fidelity}

\begin{center}
\renewcommand{\arraystretch}{1.32}
\begin{tabular}{@{}>{\raggedright\arraybackslash}p{3.3cm}
                 >{\raggedright\arraybackslash}p{5.4cm}
                 >{\raggedright\arraybackslash}p{6.5cm}@{}}
\toprule
\rowcolor{conclavebg}
{\sffamily\bfseries\color{conclaveblue}Level} &
{\sffamily\bfseries\color{conclaveblue}What it is} &
{\sffamily\bfseries\color{conclaveblue}What it can prove} \\
\midrule
\textbf{1. Constrained x86} \newline {\scriptsize what \file{exp4\_edge.py} does} &
Thread caps and a RAM ceiling on your own machine &
Memory footprint, algorithmic behaviour. Latency only as an \emph{x86} number. \\
\midrule
\rowcolor{conclaveteal!7}
\textbf{2. Emulated ARM64} \newline {\scriptsize Docker + QEMU} &
A genuine \code{aarch64} userland: ARM wheels, ARM BLAS, ARM NumPy paths &
That it \emph{installs}, \emph{runs} and \emph{fits} on ARM. Correctness.
\textbf{Not} latency. \\
\midrule
\textbf{3. Real hardware} &
An actual Pi 4, about \$50 &
Everything, including latency and energy \\
\bottomrule
\end{tabular}
\end{center}

\begin{keybox}[What emulation is genuinely for]
Level 2 answers the question that actually blocks a deployment: \emph{does the
software stack exist for this architecture at all?} Whether \code{scipy} has an
ARM wheel is a hard yes/no that emulation answers definitively, and it is a far
more common cause of failed edge deployments than being 3\,ms too slow.
\end{keybox}

\subsection{The finding that will save you a weekend}

Raspberry Pi OS ships in 32-bit (\code{armv7l}) and 64-bit (\code{aarch64})
flavours. For this dependency stack they are not equivalent:

\begin{center}
\renewcommand{\arraystretch}{1.24}
\begin{tabular}{@{}l c c@{}}
\toprule
\rowcolor{conclavebg}
{\sffamily\bfseries\color{conclaveblue}Package} &
{\sffamily\bfseries\color{conclaveblue}PyPI \code{aarch64} (64-bit)} &
{\sffamily\bfseries\color{conclaveblue}PyPI \code{armv7l} (32-bit)} \\
\midrule
numpy        & \textcolor{conclaveteal}{\textbf{wheel}} & \textcolor{conclavered}{\textbf{none}} \\
scipy        & \textcolor{conclaveteal}{\textbf{wheel}} & \textcolor{conclavered}{\textbf{none}} \\
scikit-learn & \textcolor{conclaveteal}{\textbf{wheel}} & \textcolor{conclavered}{\textbf{none}} \\
lightgbm     & \textcolor{conclaveteal}{\textbf{wheel}} & \textcolor{conclavered}{\textbf{none}} \\
psutil       & \textcolor{conclaveteal}{\textbf{wheel}} & \textcolor{conclavered}{\textbf{none}} \\
joblib       & \multicolumn{2}{c}{\textcolor{conclaveteal}{\textbf{pure Python --- fine either way}}} \\
\bottomrule
\end{tabular}
\end{center}

\begin{dangerbox}[Use the 64-bit OS]
On 32-bit Raspberry Pi OS, \code{pip install scipy} means \emph{compiling scipy
on the Pi}. That is hours of build time and it frequently dies of memory
exhaustion on a 1\,GB board.

Two ways out, in order of preference:
\begin{enumerate}[leftmargin=*,itemsep=2pt,topsep=3pt]
  \item \textbf{Install 64-bit Raspberry Pi OS} and use ordinary PyPI wheels.
        The Pi 4 and Pi Zero 2 W are both 64-bit capable. This is the right
        answer.
  \item If you are stuck on 32-bit, use \textbf{piwheels}
        (\code{piwheels.org}), which publishes pre-built \code{armv7l} wheels
        and is already configured as an index on Raspberry Pi OS. It carries
        builds for all four packages above.
\end{enumerate}
\end{dangerbox}

\subsection{Setting up emulated ARM64 on Windows}

\begin{codebox}[PowerShell --- one time]
\begin{lstlisting}[style=shell]
# 1. Docker Desktop with the WSL2 backend must be running.
docker version

# 2. Register the QEMU binfmt handlers so the kernel can execute ARM binaries.
docker run --privileged --rm tonistiigi/binfmt --install arm64

# 3. Confirm arm64 is now a build target.
docker buildx ls
\end{lstlisting}
\end{codebox}

\noindent You should see \code{linux/arm64} listed among the supported
platforms. That single \code{binfmt} step is what makes everything else work:
it teaches the Linux kernel to hand \code{aarch64} executables to QEMU
transparently.

\subsection{Build and run the Pi target}

The repository ships \file{deploy/Dockerfile.pi4} and
\file{deploy/docker-compose.pi4.yml}.

\begin{codebox}[PowerShell --- from the repo root]
\begin{lstlisting}[style=shell]
docker buildx build --platform linux/arm64 `
  -f deploy/Dockerfile.pi4 -t conclave:pi4 --load .

# sanity check: what architecture are we actually on?
docker run --rm --platform linux/arm64 conclave:pi4

# the real check, with a Pi 4's memory ceiling enforced
docker compose -f deploy/docker-compose.pi4.yml run --rm pi4 `
  python deploy/pi_check.py
\end{lstlisting}
\end{codebox}

\noindent Three details in that Dockerfile are deliberate and worth copying if
you write your own:

\begin{center}
\renewcommand{\arraystretch}{1.28}
\begin{tabular}{@{}>{\raggedright\arraybackslash}p{5.4cm}
                 >{\raggedright\arraybackslash}p{9.8cm}@{}}
\toprule
\rowcolor{conclavebg}
{\sffamily\bfseries\color{conclaveblue}Line} &
{\sffamily\bfseries\color{conclaveblue}Why} \\
\midrule
\code{--only-binary=:all:} & Makes a missing ARM wheel fail \emph{loudly} in
seconds instead of silently starting an hour-long source build. This is the
single most useful flag in the file. \\
\code{apt-get install libgomp1} & LightGBM's OpenMP runtime. The
\code{python:slim} images do not ship it, and the resulting \code{ImportError}
is cryptic. \\
\code{OMP\_NUM\_THREADS=1} & Matches the single-decision workload, for the
reason in \S\ref{sec:models}. The compose file raises it to 4 for the Pi 4
profile. \\
\bottomrule
\end{tabular}
\end{center}

\subsection{The compose profiles}

Three boards, with \textbf{real} cgroup limits --- unlike the \code{RLIMIT\_AS}
path in \file{exp4\_edge.py}, which silently does nothing on Windows:

\begin{center}
\renewcommand{\arraystretch}{1.24}
\begin{tabular}{@{}llrrl@{}}
\toprule
\rowcolor{conclavebg}
{\sffamily\bfseries\color{conclaveblue}Service} &
{\sffamily\bfseries\color{conclaveblue}Board} &
{\sffamily\bfseries\color{conclaveblue}RAM} &
{\sffamily\bfseries\color{conclaveblue}CPUs} &
{\sffamily\bfseries\color{conclaveblue}Notes} \\
\midrule
\code{pi4}    & Pi 4 Model B      & 4\,GB   & 4 & the usual target \\
\code{pi3}    & Pi 3B+            & 1\,GB   & 4 & tight but workable \\
\code{zero2w} & Pi Zero 2 W       & 512\,MB & 1 & the hard case \\
\bottomrule
\end{tabular}
\end{center}

\begin{keybox}[\code{memswap\_limit} equal to \code{mem\_limit} disables swap]
Without that, Docker lets the container swap and your ``512\,MB'' test quietly
becomes a 1\,GB test. Setting them equal is what makes the ceiling real. This is
the most commonly botched part of container-based memory testing.
\end{keybox}

\subsection{What \code{pi\_check.py} reports, and why it is split}

\file{deploy/pi\_check.py} trains a small agent, exercises it, and labels every
line \code{[valid]} or invalid, detecting emulation by looking for a real Pi
signature in \file{/proc/cpuinfo} and \file{/proc/device-tree/model}:

\begin{codebox}[One captured run on native 2-core x86 --- your numbers will differ]
\begin{lstlisting}[style=out]
  architecture   : x86_64
  visible cores  : 2
  RAM ceiling    : 3915 MB
  emulated       : no
====================================================================
  [valid]  agent artefact      :    186.2 KB
  [valid]  training accuracy   :     98.19 %
  [valid]  peak RSS            :    194.9 MB
  [valid]  fits RAM ceiling    : YES
  [valid]  training time       :       2.0 s

  [valid (native)]
           median decision     :      2.41 ms
           p99 decision        :      2.92 ms
\end{lstlisting}
\end{codebox}

\noindent Run the same script under the \code{arm64} container and the header
flips to \code{emulated: YES}, the timing block is marked invalid, and the
script tells you so in plain language.

\begin{warnbox}[Latency is noisy; do not quote a single run]
Repeated runs of the same benchmark on the same machine vary by tens of
percent, depending on thermal state and what else is running. Across three runs
of \file{exp4\_edge.py} on the same container the host median moved between
2.4 and 2.8\,ms. If a latency number is going into a paper, report the median
of several independent runs and say how many --- and re-run every experiment
that feeds a table before you submit, because a stale table is the easiest
mistake to make and the hardest to spot.
\end{warnbox}

\subsection{What you may and may not claim}

\begin{center}
\renewcommand{\arraystretch}{1.3}
\begin{tabular}{@{}>{\raggedright\arraybackslash}p{5.0cm}
                 >{\raggedright\arraybackslash}p{10.2cm}@{}}
\toprule
\rowcolor{conclavebg}
{\sffamily\bfseries\color{conclaveblue}Measurement} &
{\sffamily\bfseries\color{conclaveblue}Status under emulation} \\
\midrule
Model / artefact size & \textcolor{conclaveteal}{\textbf{Valid.}} Architecture-independent bytes. \\
Peak RSS & \textcolor{conclaveteal}{\textbf{Valid.}} The cgroup limit is real, and pointer width is genuinely 64-bit ARM. \\
Fits in RAM (yes/no) & \textcolor{conclaveteal}{\textbf{Valid}}, and this is the claim that most often decides a deployment. \\
Correctness / accuracy & \textcolor{conclaveteal}{\textbf{Valid.}} Same wheels, same code paths. \\
Wheel availability & \textcolor{conclaveteal}{\textbf{Valid}}, and only obtainable this way. \\
\midrule
Decision latency & \textcolor{conclavered}{\textbf{Invalid.}} QEMU is typically 5--20$\times$ slower and not uniformly so. \\
Energy per decision & \textcolor{conclavered}{\textbf{Invalid.}} Derived from latency. \\
Throughput & \textcolor{conclavered}{\textbf{Invalid.}} \\
\bottomrule
\end{tabular}
\end{center}

\begin{warnbox}[How to phrase this in the paper]
The paper currently reports estimated on-device latency scaled by a relative
single-core throughput ratio, and labels it an estimate. That is defensible.
What is \emph{not} defensible is presenting a QEMU number as a Pi measurement.
If a reviewer asks for hardware numbers, the honest answer is: ``memory and
correctness were verified on an emulated \code{aarch64} target with enforced
cgroup limits; latency is a stated scaling model, and we have not measured on
silicon.'' Buying a Pi 4 is the cheapest way to close that gap.
\end{warnbox}

\subsection{When you do get real hardware}

\begin{codebox}[On the Pi, over SSH --- 64-bit Raspberry Pi OS]
\begin{lstlisting}[style=shell]
uname -m                       # want: aarch64
sudo apt update && sudo apt install -y python3-venv python3-dev libgomp1

python3 -m venv ~/conclave-env
source ~/conclave-env/bin/activate

git clone https://github.com/<user>/conclave && cd conclave
pip install --only-binary=:all: -e ".[boost]"

python deploy/pi_check.py      # now every line is valid
pytest -q                      # 9 smoke tests; paper-claim tests need results/
\end{lstlisting}
\end{codebox}

\noindent Two practical notes for real hardware:

\begin{itemize}[leftmargin=*,itemsep=3pt]
  \item \textbf{Thermal throttling is real.} A Pi 4 under sustained load drops
        clock speed. Run \code{vcgencmd measure\_temp} alongside your benchmark
        and discard runs above about 80\,\textdegree C, or report that you used
        a heatsink.
  \item \textbf{Pin the governor} before timing anything:
        \code{sudo cpufreq-set -g performance}. Otherwise the first few
        decisions are measured at idle clock and your p99 is meaningless.
\end{itemize}

\subsection{Deploying a trained agent, not the training code}

For a field device you ship the artefact, not the pipeline:

\begin{codebox}[On your laptop]
\begin{lstlisting}[style=py]
import joblib
joblib.dump({"fz": fz, "brain": brain, "gate": brain.gate},
            "agent.joblib", compress=3)     # ~200 KB
\end{lstlisting}
\end{codebox}

\begin{codebox}[On the Pi --- no pandas, no matplotlib, no training data]
\begin{lstlisting}[style=py]
import joblib
from conclave import Agent
from conclave.core.types import Observation

b = joblib.load("agent.joblib")
agent = Agent(brain=b["brain"], featurizer=b["fz"],
              registry=registry, gate=b["gate"])

while True:
    obs = Observation(features=read_sensors(), text=read_radio())
    d = agent.decide(obs)
    log(d.action, d.confidence, d.disagreement, d.gate)
\end{lstlisting}
\end{codebox}

\noindent The runtime dependency set collapses to \code{numpy}, \code{scipy},
\code{scikit-learn}, \code{joblib} --- and \code{lightgbm} only if you use
\code{ToolRouter}. That is roughly 60\,MB of wheels and a 200\,KB agent.
